% Opening of the old chapter c24, its \chapter command removed; generated
\label{chap:hitting}\label{ch:ht}

The reciprocal of the normalised $\xi$ function is the Laplace transform of a probability law, the law of the clock $T$ of Section~\ref{ch:clock}. Under the Riemann Hypothesis that law is the hitting time of a Krein string and the explosion time of a pure birth chain, and its density is Kent's eigenfunction expansion over the zeros. We ask whether the argument can be run backwards, so that RH follows from the limiting behaviour of hitting times, and show where the content of such an argument must lie. We test approximating laws built with no assumption on the zeros (Jensen laws, Gram smoothed zeros, truncated Euler products, boundaries of Brownian motion), prove a universal stable limit for the prime side hitting times and a rigidity theorem for factorisations of the clock law, and study the density of the clock directly: its residue expansion, its log concavity and the failure of hyperbolic complete monotonicity. No step proves RH.

The material is based on \cite{P26ht,PolsonResidue26,PolsonEuler26}; its numbers were recomputed for this book, and results no independent computation reached are marked unchecked. Throughout, $G(u)=\xi(\tfrac12+\sqrt u)/\xi(\tfrac12)$ as in Section~\ref{ch:clock}, the clock $T$ has law $\mu$ and Laplace transform $1/G$, the canonical function is the heat trace $W$ (Chapter~\ref{chap:heat}), and the de Bruijn and Newman kernel is $\Phi_{\mathrm{dBN}}$ of Section~\ref{ch:dbn}, not the kernel $\Phi$ of \eqref{eq:ch2:Phi}. Results proved elsewhere are cited, not restated: the self decomposability of $\mu$ and the equivalence of the GGC property with RH (Theorem~\ref{thm:ch12:idsd}, Proposition~\ref{prop:rigid}), the spectral expansion and the Kent approximants (Section~\ref{ch:clocklaw}), and the closure of the GGC class under weak limits (Corollary~\ref{cor:weak:closure}).
